% GENERATED FILE. Edit canonical lessons, then run npm run generate:books.
% canonical-source-hash: fnv1a64:aa484fe0b380ff4b
% canonical-lessons: JA-C274-jishin, JA-C274-kuzureru, JA-C274-korogaru, JA-C274-osoroshii, JA-C274-hinan

\chapter{When the Earth Shook}
\label{ch:ja-when-the-earth-shook}
\hlchaptermodality{274}

\begin{chapteropening}
\textbf{By the end of this chapter:} \emph{I can say what an earthquake did to the town.}

Complete the last lesson of chapter 274: I can say what an earthquake did to the town.
\end{chapteropening}

\section[\texorpdfstring{\ja{じしん} (jishin)}{じしん (jishin)}]{\ja{じしん} (jishin) — an earthquake}

\label{lesson:JA-C274-jishin}



\begin{quote}
Before the new one: say the Japanese for to sue, then the Japanese for evidence.
\end{quote}

\subsection*{You'll want to know: \ja{じしん}}
\textbf{\ja{じしん}} — \emph{jishin} — \textquotedblleft{}an earthquake\textquotedblright{}.

\begin{quote}
\textbf{\ja{きのう} \ja{じしんが} \ja{ありました。}} — \emph{kinō jishin ga arimashita} — \textquotedblleft{}There was an earthquake yesterday.\textquotedblright{}
\end{quote}

\begin{cousinweb}
\textbf{\ja{じしん}} (\emph{jishin}) is Sino-Japanese (kango): two Chinese characters read with their on'yomi, readings borrowed from Chinese. They are \emph{ji}, \textquotedblleft{}ground\textquotedblright{} (Middle Chinese \emph{dijH}) and \emph{shin}, \textquotedblleft{}quake\textquotedblright{} (\emph{tsyinH}) (KANJIDIC2; Middle Chinese in Baxter's transcription).
\end{cousinweb}

\begin{grammarlens}[title={What you've built}]
One more word for telling what happened.
\end{grammarlens}

\subsection*{Your turn}
\begin{itemize}
\raggedright
  \item \emph{Say it:} \emph{jishin}
  \item \emph{Say it:} \emph{jishin}, once more
  \item \emph{Say it:} \emph{kinō jishin ga arimashita} — There was an earthquake yesterday.
  \item \emph{Recall:} say the Japanese for to sue, then the Japanese for evidence, then say \emph{jishin} again
\end{itemize}

\begin{tcolorbox}[breakable,colback=teal!4,colframe=teal!35!black,title={Before you move on}]
What does \ja{じしん} mean? (\textquotedblleft{}An earthquake\textquotedblright{}.) Say it once more.
\end{tcolorbox}
\section[\texorpdfstring{\ja{くずれる} (kuzureru)}{くずれる (kuzureru)}]{\ja{くずれる} (kuzureru) — to crumble; to give way}

\label{lesson:JA-C274-kuzureru}



\begin{quote}
Before the new one: say the Japanese for evidence, then the Japanese for an earthquake.
\end{quote}

\subsection*{You'll want to know: \ja{くずれる}}
\textbf{\ja{くずれる}} — \emph{kuzureru} — \textquotedblleft{}to crumble; to give way\textquotedblright{}.

\begin{quote}
\textbf{\ja{じしんで} \ja{ふるい} \ja{たてものが} \ja{くずれました。}} — \emph{jishin de furui tatemono ga kuzuremashita} — \textquotedblleft{}Old buildings came down in the earthquake.\textquotedblright{}
\end{quote}

Polite \textbf{\ja{くずれます}}, \emph{kuzuremasu}; an ichidan verb. Past (polite): \textbf{\ja{やまが} \ja{くずれました。}}, \emph{yama ga kuzuremashita}.

\begin{cousinweb}
\textbf{\ja{くずれる}} (\emph{kuzureru}) is a native Japanese word (wago), inherited rather than borrowed from Chinese; its kanji was attached later as a kun'yomi (JMdict).
\end{cousinweb}

\begin{grammarlens}[title={What you've built}]
One more word for telling what happened.
\end{grammarlens}

\subsection*{Your turn}
\begin{itemize}
\raggedright
  \item \emph{Say it:} \emph{kuzureru}
  \item \emph{Say it:} \emph{kuzureru}, once more
  \item \emph{Say it:} \emph{jishin de furui tatemono ga kuzuremashita} — Old buildings came down in the earthquake.
  \item \emph{Recall:} say the Japanese for evidence, then the Japanese for an earthquake, then say \emph{kuzureru} again
\end{itemize}

\begin{tcolorbox}[breakable,colback=teal!4,colframe=teal!35!black,title={Before you move on}]
What does \ja{くずれる} mean? (\textquotedblleft{}To crumble; to give way\textquotedblright{}.) Say it once more.
\end{tcolorbox}
\section[\texorpdfstring{\ja{ころがる} (korogaru)}{ころがる (korogaru)}]{\ja{ころがる} (korogaru) — to roll; to lie about}

\label{lesson:JA-C274-korogaru}



\begin{quote}
Before the new one: say the Japanese for an earthquake, then the Japanese for to crumble.
\end{quote}

\subsection*{You'll want to know: \ja{ころがる}}
\textbf{\ja{ころがる}} — \emph{korogaru} — \textquotedblleft{}to roll; to lie about\textquotedblright{}.

\begin{quote}
\textbf{\ja{いしが} \ja{みちに} \ja{ころがっていました。}} — \emph{ishi ga michi ni korogatte imashita} — \textquotedblleft{}Stones were lying all over the road.\textquotedblright{}
\end{quote}

Polite \textbf{\ja{ころがります}}, \emph{korogarimasu}; a godan verb. Past (polite): \textbf{\ja{いしが} \ja{ころがりました。}}, \emph{ishi ga korogarimashita}.

\begin{cousinweb}
\textbf{\ja{ころがる}} (\emph{korogaru}) is a native Japanese word (wago), inherited rather than borrowed from Chinese; its kanji was attached later as a kun'yomi (JMdict).
\end{cousinweb}

\begin{grammarlens}[title={What you've built}]
One more word for telling what happened.
\end{grammarlens}

\subsection*{Your turn}
\begin{itemize}
\raggedright
  \item \emph{Say it:} \emph{korogaru}
  \item \emph{Say it:} \emph{korogaru}, once more
  \item \emph{Say it:} \emph{ishi ga michi ni korogatte imashita} — Stones were lying all over the road.
  \item \emph{Recall:} say the Japanese for an earthquake, then the Japanese for to crumble, then say \emph{korogaru} again
\end{itemize}

\begin{tcolorbox}[breakable,colback=teal!4,colframe=teal!35!black,title={Before you move on}]
What does \ja{ころがる} mean? (\textquotedblleft{}To roll; to lie about\textquotedblright{}.) Say it once more.
\end{tcolorbox}
\section[\texorpdfstring{\ja{おそろしい} (osoroshii)}{おそろしい (osoroshii)}]{\ja{おそろしい} (osoroshii) — frightening; dreadful}

\label{lesson:JA-C274-osoroshii}



\begin{quote}
Before the new one: say the Japanese for to crumble, then the Japanese for to roll.
\end{quote}

\subsection*{You'll want to know: \ja{おそろしい}}
\textbf{\ja{おそろしい}} — \emph{osoroshii} — \textquotedblleft{}frightening; dreadful\textquotedblright{}.

\begin{quote}
\textbf{\ja{それは} \ja{おそろしい} \ja{じしんでした。}} — \emph{sore wa osoroshii jishin deshita} — \textquotedblleft{}It was a frightening earthquake.\textquotedblright{}
\end{quote}

An i-adjective.

\begin{cousinweb}
\textbf{\ja{おそろしい}} (\emph{osoroshii}) is a native Japanese word (wago), inherited rather than borrowed from Chinese; its kanji was attached later as a kun'yomi (JMdict).
\end{cousinweb}

\begin{grammarlens}[title={What you've built}]
One more word for telling what happened.
\end{grammarlens}

\subsection*{Your turn}
\begin{itemize}
\raggedright
  \item \emph{Say it:} \emph{osoroshii}
  \item \emph{Say it:} \emph{osoroshii}, once more
  \item \emph{Say it:} \emph{sore wa osoroshii jishin deshita} — It was a frightening earthquake.
  \item \emph{Recall:} say the Japanese for to crumble, then the Japanese for to roll, then say \emph{osoroshii} again
\end{itemize}

\begin{tcolorbox}[breakable,colback=teal!4,colframe=teal!35!black,title={Before you move on}]
What does \ja{おそろしい} mean? (\textquotedblleft{}Frightening; dreadful\textquotedblright{}.) Say it once more.
\end{tcolorbox}
\section[\texorpdfstring{\ja{ひなん} (hinan)}{ひなん (hinan)}]{\ja{ひなん} (hinan) — evacuation; taking refuge}

\label{lesson:JA-C274-hinan}



\begin{quote}
Before the new one: say the Japanese for to roll, then the Japanese for frightening.
\end{quote}

\subsection*{You'll want to know: \ja{ひなん}}
\textbf{\ja{ひなん}} — \emph{hinan} — \textquotedblleft{}evacuation; taking refuge\textquotedblright{}.

\begin{quote}
\textbf{\ja{じしんの} \ja{あと、ちちが} \ja{ひなんを} \ja{言}\ja{い}\ja{出}\ja{しました。}} — \emph{jishin no ato, chichi ga hinan o iidashimashita} — \textquotedblleft{}After the earthquake, my father said we should leave.\textquotedblright{}
\end{quote}

\begin{cousinweb}
\textbf{\ja{ひなん}} (\emph{hinan}) is Sino-Japanese (kango): two Chinese characters read with their on'yomi, readings borrowed from Chinese. They are \emph{hi}, \textquotedblleft{}avoid\textquotedblright{} (Middle Chinese \emph{bjieH}) and \emph{nan}, \textquotedblleft{}trouble\textquotedblright{} (\emph{nanH}) (KANJIDIC2; Middle Chinese in Baxter's transcription).
\end{cousinweb}

\begin{grammarlens}[title={What you've built}]
One more word for telling what happened.
\end{grammarlens}

\subsection*{Your turn}
\begin{itemize}
\raggedright
  \item \emph{Say it:} \emph{hinan}
  \item \emph{Say it:} \emph{hinan}, once more
  \item \emph{Say it:} \emph{jishin no ato, chichi ga hinan o iidashimashita} — After the earthquake, my father said we should leave.
  \item \emph{Recall:} say the Japanese for to roll, then the Japanese for frightening, then say \emph{hinan} again
\end{itemize}

\begin{tcolorbox}[breakable,colback=teal!4,colframe=teal!35!black,title={Before you move on}]
What does \ja{ひなん} mean? (\textquotedblleft{}Evacuation; taking refuge\textquotedblright{}.) Say it once more.
\end{tcolorbox}
